
For the positive cone, $H^2_{\rm glob}\to H^2(\mathbf R,T)$ is onto:
any real component's image in $H^0(\Q,D)^\vee$ can be canceled by
a component at two, since restriction
$H^0(\Q,D)\to H^0(\Q_2,D)$ is injective. Thus its exact sequences are
\[
0\to T^+\xrightarrow{a}P\to K\to0,\qquad
0\to\operatorname{coker}(H^1_{\rm glob}\to H^1_{\mathbf R})
 \to H^2(C_h^+)\to H^2_{\rm glob}\to H^2_{\mathbf R}\to0.
\]
If $k$ and $p_0$ are the torsion lengths of $K$ and $P$, $m$ is the
free index of $K$ in global $H^1$, and $a_s$ is the saturation index
of the real line in $P/P_{\rm tors}$, these sequences give
\[
\ell\operatorname{coker}(H^1_{\rm glob}\to H^1_{\mathbf R})
 =r_1-g+k-m,\qquad k=p_0+a_s.
\]
The wedge $2^{-c_E}a\wedge b$ contributes
$a_s+1+v_{\rm free}(z)-m$, and the inverse determinant comparison
contributes $p_0-\ell H^2(C_h^+)$. Substitution leaves
$1+v_{\rm free}(z)+i-s(h)-\sum_vg_v+2g-r_1$.
Finally $v_2\exp^*z=v_{\rm free}(z)+i+e_2$, proving (E.7)
without specializing a primitivity assertion from the punctured
family.
Apply (K3) to the retained class $z=F_E(Z_F)_h$ used in the wedge
above. The primitive-conductor trace and the invariant
$\sqrt h\,e_{\chi_h}$ basis have already been computed in
paragraph~6, including discriminant $4h$. Absorb the fixed
smoothing and old-support factors into $\epsilon_{\rm rec}$.
Then
\[
 v_2\exp^*z
 =v_2L_*(h)+\sum_{q\ {\rm unused}}v_2\#E^h(\mathbf F_q)
                         +\epsilon_{\rm rec},\qquad
 |\epsilon_{\rm rec}|\le C_E.
\]
At unused good primes $g_q=v_2\#E^h(\mathbf F_q)$, so these
terms cancel. At an active good prime the inertia invariants in
$E^h[2^\infty]$ are $E[2]$, of length two for the present curve.
The remaining local types have fixed bounds, including global
torsion by its injection at two.
Thus
\[
 v_2u_h^{\rm ev}(0)=v_2L_*(h)-2\omega(h)-s(h)+\epsilon_h,\qquad
 |\epsilon_h|\le C_{\rm loc}.                                  \tag{E.8}
\]
The central rational positive complex has only rank-two $H^1$, so
this is specialization of (E.4). Corank-zero duality also shows
$u_h^{\rm ev}(0)=0$ if and only if $L_*(h)=0$.
At $L_*(h)\ne0$, Lemma~\ref{lem:forward-rank} supplies the
corank-zero and finite-$\Sha$ hypotheses just used.
Equations (E.5) and (E.8), already with $G=1$, yield
\[
          v_2L_*(h)\ge2\omega(h)+s(h)-c_E-C_{\rm loc}.             \tag{E.9}
\]
No binary-family dimension is needed for this lower bound.
\end{proof}

